\documentclass[]{article}

\usepackage{amsmath}
\usepackage{multirow}
\usepackage{natbib}
\usepackage{graphicx} 


\begin{document}

The transport of particles and scalars by turbulent flows is a fundamental process in physics, governing phenomena from pollutant dispersion in the atmosphere to mixing in astrophysical plasmas. While classical diffusion theory provides a robust framework for random transport, its core assumption of short-correlated, Gaussian fluctuations is frequently violated in turbulence. Turbulent velocity fields are characterized by long-range correlations and intermittent, non-Gaussian statistics, which often give rise to anomalous transport regimes where particle dispersion deviates significantly from classical predictions. This departure necessitates more general theoretical frameworks capable of capturing the complex, multi-scale nature of turbulent diffusion.

A particularly suitable framework for describing such anomalous transport is provided by fractional calculus. In this approach, the evolution of the probability density function of a tracer particle, $P(x,t)$, is modeled by a fractional diffusion equation, $\partial_t P(x,t) = -D_\alpha(-\Delta)^{\alpha/2} P(x,t)$. The Lévy exponent $\alpha$ in this equation characterizes the transport dynamics; the value $\alpha=2$ recovers the standard diffusion equation describing Brownian motion, while values of $\alpha < 2$ correspond to superdiffusive processes known as Lévy flights, which are driven by jump distributions with heavy tails. This formulation establishes a direct mathematical link between the macroscopic transport behavior and the statistical properties of the underlying particle displacements.

A key theoretical prediction, derived from Renormalization Group (RG) arguments, connects the Lagrangian transport exponent $\alpha$ to the statistical geometry of the Eulerian velocity field. Specifically, the theory posits that in the asymptotic limit of long times and large scales, the system flows to a fixed point where the Lévy exponent is determined by the velocity field's spectral roughness $\xi$ via the relation $\alpha = 2/\xi$. The roughness exponent $\xi$ is, in turn, directly related to the second-order structure function of the velocity field, which quantifies its spatial regularity. This relationship offers a powerful and predictive bridge between the static, Eulerian statistics of the fluid and the dynamic, path-dependent Lagrangian statistics of tracer particles.

However, this theoretical link describes an asymptotic state, and its observability in any real or simulated system with finite scales and finite duration is not guaranteed \citep{onizawa2026finitetimeobservabilityoscillatoryinstabilities}. The central problem this study addresses is to determine the conditions under which this predicted fractional regime is physically realized and to identify the mechanisms that may hinder its emergence. It remains unclear how, and over what timescales, a system transitions from its initial short-time behavior, which is typically ballistic or diffusive, to the predicted anomalous fixed point \citep{onizawa2026finitetimeobservabilityoscillatoryinstabilities}. The practical relevance of the theory hinges on understanding this crossover.

In this paper, we systematically investigate the emergence of these asymptotic fractional dynamics \citep{edelman2022cyclesasymptoticallystablechaotic} across a hierarchy of numerical turbulence models, including stochastic multifractal cascades, the synthetic Kraichnan model, and the deterministic Lorenz-96 system. Our approach moves beyond a simple measurement of global scaling exponents by explicitly tracking the RG flow of an effective, time-dependent exponent, $\alpha(\tau)$. This method allows us to directly visualize the crossover from the initial transport regime toward the predicted asymptotic state \citep{majdoub2025longtimeasymptoticssubordinatedfractional}. By doing so, we aim to delineate the critical factors—such as sufficient scale separation, simulation duration, and spatial dimensionality—that govern whether the fractional operator defined by the Eulerian roughness provides a physically manifest description of turbulent transport or remains a purely theoretical, unobservable limit.

\end{document}
                